
% 05. The Continuum of Life and Death: The Spectrum of Dynamicity.tex

\section{생명과 죽음의 연속선: 역동성(Dynamicity)의 스펙트럼}
우리는 세상을 '살아 있는 것'과 '죽어 있는 것'으로 엄격하게 구분하는 데 너무도 익숙하다. 흔히 돌과 강물은 죽어 있는 무생물이며, 나무와 물고기, 그리고 인간은 살아 숨 쉬는 생명체라고 철석같이 믿는다. 생명은 우주의 일반적인 무기물들과는 완전히 다른 차원의 원리로 움직이는 특별하고도 고결한 기적이라 찬미해 온 것이다. 그래서 우리는 생명과 죽음을 서로 완벽하게 단절된 이분법적 대립항으로 설정하고, 그 사이에 결코 건널 수 없는 뚜렷한 경계선이 존재한다고 여긴다. 그러나 디스턴스(DISTANCE)의 서늘한 렌즈를 들이대고 우주 전체를 하나의 거대한 모멘텀 해소의 흐름으로 바라본다면, 이 견고한 이분법은 완벽한 환상임이 폭로된다.

시공간의 바탕화면을 치워버린 우주의 실제 하드웨어 위에는 수많은 '에너지 아일랜드(Energy Island)'들이 흩어져 있다. 별도 하나의 거대한 아일랜드이고, 행성도 하나의 아일랜드이며, 바위와 대기권, 그리고 인간의 세포 역시 각자의 모멘텀을 품고 있는 에너지 아일랜드들이다. 우주에서 벌어지는 모든 현상은 오직 이 아일랜드들 사이에 존재하는 '에너지 갭(낙차)'을 해소하기 위해 맹렬하게 교류하는 과정에서 파생된다. 즉, 우주의 진정한 본질은 생성과 소멸이라는 극적인 마법이 아니라, 수많은 에너지 아일랜드들이 모멘텀을 주고받으며 결합하고 해체하는 영구적인 '이합집산(Gathering and Scattering)'일 뿐이다. 이 장엄한 이합집산의 궤적 위에서, 살아 있다는 것과 죽어 있다는 것은 결코 서로 다른 두 개의 세계가 아니다. 그것은 선형적인 에너지 스케일 위에서 나타나는 '역동성(Dynamicity)'의 차이를 나타내는 연속적인 상태점들일 뿐이다.

우리가 흔히 '죽어 있다'고 표현하는 바위나 쇳덩어리는 에너지의 잠재성(포텐셜)이 극도로 적게 묶여 있어 다른 존재와의 에너지 흐름이나 교환이 거의 멈춰버린 '다이내믹시티(역동성)가 굉장히 낮은 엔터티'일 뿐이다. 반면 우리가 '살아 있다'고 부르는 존재들은 포텐셜 에너지가 높아 에너지 교류를 촉진하려는 모멘텀이 극도로 강한, 즉 '다이내믹시티가 굉장히 높은 엔터티'에 불과하다. 결국 생명과 무생물 사이의 차이는 본질적인 질적 차이가 아니라, 우주의 에너지를 섞어내고 교류하는 역동성(모멘텀의 강도)의 양적 차이에 지나지 않는다. 우리는 단지 우주의 무한한 이합집산 과정 중 특정한 역동성의 구간을 자의적으로 잘라내어 '생명'이라 부르고, 그 역동성이 현저히 낮아진 구간을 '죽음'이라 부르며 착각하고 있었던 것이다.

이러한 통찰은 기존 생물학이 생명을 대하던 낡은 태도를 완전히 전복시킨다. 생명은 우주의 법칙(엔트로피 증가)을 거스르며 무질서의 바다 위에서 기적적으로 스스로의 질서를 구축해 낸 예외적 반역자가 결코 아니다. 오히려 생명은 가장 극단적인 모멘텀을 품고, 우주의 에너지 갭을 누구보다 맹렬하게 깎아내려 평활화(Equalization)시키기 위해 형성된 가장 정교한 '로컬 모멘텀 구조(Local Momentum Structure)'이다.
생명체가 먹고, 소화하고, 호흡하고, 배설하는 모든 과정은 무엇인가? 그것은 우주에서 고도로 압축된 포텐셜 에너지(식량, 태양빛)를 거칠게 빨아들여, 자신의 작은 내부 질서를 유지한다는 핑계로 주변 환경에 훨씬 거대하고 무질서한 열과 폐기물(저품질 에너지)을 폭발적으로 흩뿌리는 행위이다. 생명은 내부에 에너지를 가만히 쥐고 있기 때문에 살아 있는 것이 아니다. 외부와 쉴 새 없이 에너지를 교환하고 분산시키기 때문에, 즉 '우주의 평활화를 지연시키는 것이 아니라 극단적으로 가속하고 있기 때문에' 살아 있는 것이다.

그렇다면, 이 찬란한 생명이 맞이하는 '죽음'의 진짜 실체는 무엇인가? 죽음은 단순히 심장의 박동이 멈추거나 존재가 우주에서 삭제되는 종료 버튼이 아니다. 죽음은 더 큰 에너지 흐름(생태계와 우주)과의 맹렬했던 교류가 끊어지는 '절대적 단절(디스턴스)'의 사건이다. 살아 있는 동안 생명체는 외부의 에너지를 받아들이고 섞어서 내보내는 맹렬한 소용돌이(흐름 속의 매듭)였으나, 그 교류가 멈추는 순간 생명체는 완벽히 고립된 '에너지 아일랜드'로 전락한다. 내부에 남아 있던 모멘텀이 소진되고 다이내믹시티(역동성)가 극도로 낮아지며, 국소적 차원에서의 작은 '열적 평형(죽음)'에 도달하는 것이다. 그러나 우주의 거시적 관점에서는 어떤 것도 끝나지 않았다. 죽은 생명체를 구성하던 물질과 에너지는 분해자(미생물)에 의해 해체되어 토양과 대기, 그리고 다른 생명체의 일부로 다시 맹렬하게 연결된다. 죽음은 종료가 아니라 우주의 평활화라는 웅장한 극장 안에서 '새로운 연결과 이합집산을 위한 궤도 전환'일 뿐이다. 

결국 이 우주를 가득 채우고 있는 것은 살아 숨 쉬는 것과 죽어 있는 것의 대립이 아니다. 단지 우주 전체가 완전한 평활화(엔트로피 최대화)를 향해 끝없이 추락하는 과정 속에서, 모멘텀을 맹렬히 해소하며 요동치는 높은 역동성의 덩어리(생명)들과, 그 역동성을 잃고 흩어지는 덩어리(무생물)들이 빽빽하게 나열된 '하나의 완벽한 연속선(Continuum)'만이 존재할 뿐이다. 인간을 포함한 모든 생명은 우주의 텅 빈 무대 위에 내던져진 외로운 이방인이 아니다. 우리는 138억 년의 거대한 우주적 추락(평활화) 속에서, 우주가 스스로의 모멘텀을 맹렬히 섞고 분산시키기 위해 잠시 피워낸 가장 뜨겁고 경이로운 '에너지 소용돌이'인 것이다.

